\mode*

\section{En klass för bråk}

I förra kapitlet lärde vi oss grunderna i objektorienterad programmering med 
klasser och objekt. Nu ska vi fördjupa oss i mer avancerade tekniker genom 
att bygga en egen klass för att representera bråk (matematiska bråktal som 
\(\frac{1}{2}\) eller \(\frac{3}{4}\)).

Detta är ett utmärkt exempel eftersom:
\begin{itemize}
  \item Bråk har en naturligt sammansatt struktur (täljare och nämnare).
  \item Vi behöver implementera flertalet aritmetiska operationer och 
    jämförelser.
  \item Vi kan se hur dundermetoder används i praktiken.
  \item Vi lär oss om egenskaper (\foreignlanguage{english}{properties})
\end{itemize}

\subsection{Konstrueraren}

Vi börjar med att definiera hur vår \mintinline{python}{Fraction}-klass ska 
skapas. Konstrueraren (\mintinline{python}{__init__}-metoden) är den metod 
som anropas när vi skapar ett nytt objekt.

\begin{exercise}
  Vilka argument behöver vi för att skapa ett bråk? Tänk på vad ett bråk 
  består av matematiskt.
  Vad är användbara argument att skapa ett nytt bråk ifrån?
\end{exercise}

Ett bråk består av två delar: täljare och nämnare.
Men vad händer om vi vill skapa ett heltal som bråk?
Eller skapa ett bråk från ett annat bråk?
Konstrueraren kan göras ganska flexibel.
Den kan hantera flera olika sätt att skapa bråk:
\begin{enumerate}
  \item \mintinline{python}{Fraction(1, 2)}: Standard täljare och nämnare.
  \item \mintinline{python}{Fraction(2)}: Heltal som bråk (nämnare blir 1).
  \item \mintinline{python}{Fraction(Fraction(1, 2))}: Kopiera ett bråk.
  \item \mintinline{python}{Fraction(Fraction(1, 2), 2)}: Dela ett bråk.
\end{enumerate}

\begin{exercise}
  Hur tror du att konstrueraren vet vilken typ av argument den får?
  Hur kan den hantera både heltal och \mintinline{python}{Fraction}-objekt som 
  första argument?
\end{exercise}

\begin{frame}[fragile]
  \inputminted[linenos,firstline=3,lastline=13]{python}{examples/myfrac_base.py}

  \begin{example}
    Några sätt att skapa bråk:
    \begin{minted}{python}
      frac = Fraction(1, 2)
      int_as_frac = Fraction(2)
      frac_from_frac = Fraction(Fraction(1, 2))
      half_frac = Fraction(Fraction(1, 2), 2)
    \end{minted}
  \end{example}
\end{frame}

Konstrueraren använder \mintinline{python}{isinstance()} för att kontrollera 
typen på argumentet. Detta är ett exempel på \emph{polymorfism}---samma metod 
kan hantera olika typer av indata.


\subsection{Attribut som egenskaper}

I förra kapitlet såg vi hur vi använde \enquote{\foreignlanguage{english}{getter}}-metoder 
som \mintinline{python}{get_name()} för att läsa attribut. Nu ska vi lära oss 
ett mer elegant sätt: egenskaper (\foreignlanguage{english}{properties}).

\begin{exercise}
  Vad är skillnaden mellan att skriva \mintinline{python}{frac.get_nominator()} 
  och \mintinline{python}{frac.nominator}? Vilket känns mer naturligt?
\end{exercise}

Egenskaper låter oss använda attribut-liknande syntax medan vi fortfarande 
behåller kontrollen över hur data läses och skrivs.

\begin{frame}[fragile]
  \inputminted[linenos,firstline=3,lastline=3]{python}{examples/myfrac_base.py}
  \dots
  \inputminted[linenos,autogobble=false,firstline=15,lastline=24]{python}{examples/myfrac_base.py}

  \begin{example}
    Hur vi använder egenskaperna:
    \begin{minted}{python}
      frac = Fraction(1, 2)
      print(f"{frac.nominator}/{frac.denominator}")
    \end{minted}
  \end{example}
\end{frame}

Här använder vi \mintinline{python}{@property}-dekoratorn för att göra 
metoderna \mintinline{python}{nominator()} och \mintinline{python}{denominator()} 
tillgängliga som om de vore attribut. Detta ger en renare syntax.

\begin{exercise}
  Varför tror du det är bra att använda egenskaper istället för att göra 
  attributen publika?
\end{exercise}

Egenskaper ger oss inkapsling som vi får via getters och setters, men samtidigt 
ser det ut som direkt attributåtkomst.
Vi kan ändra den interna implementationen utan att ändra hur användare av 
klassen interagerar med den.
Vi kan också lägga till validering eller beräkningar när värden läses eller 
skrivs.

Låt oss jämföra två sätt att arbeta med attribut:

\begin{frame}[fragile]
  \begin{flushleft}
  \begin{minipage}[t]{0.49\linewidth}
    \begin{minted}{python}
        class Class:
          def __init__(self, attribute):
            self.__attribute = attribute

          @property
          def attribute(self):
            return self.__attribute

          @property.setter
          def attribute(self, value):
            self.__attribute = value

        obj = Class()
        print(obj.attribute)
        obj.attribute = new_value
    \end{minted}
  \end{minipage}\hfill
  \begin{minipage}[t]{0.49\linewidth}
    \begin{minted}{python}
        class Class:
          def __init__(self, attribute):
            self.__attribute = attribute

          def get_attribute(self):
            return self.__attribute

          def set_attribute(self, value):
            self.__attribute = value

        obj = Class()
        print(obj.get_attribute())
        obj.set_attribute(new_value)
    \end{minted}
  \end{minipage}
  \end{flushleft}
\end{frame}

I vänster kolumn ser vi hur vi kan definiera både en \foreignlanguage{english}{getter} 
och en \foreignlanguage{english}{setter} med \mintinline{python}{@property}. 
I höger kolumn ser vi det traditionella sättet med explicita metoder.


\subsection{Typkonvertering}

Nu har vi ett bråk som kan skapas och vars täljare och nämnare kan läsas. Men 
hur ska det visas för användaren? Och hur konverterar vi det till andra typer?

\begin{exercise}
  Vad tror du händer om vi försöker skriva ut ett \mintinline{python}{Fraction}-objekt 
  med \mintinline{python}{print()}? Hur ska det se ut?
\end{exercise}

Vi behöver implementera dundermetoder för att konvertera vårt bråk till 
olika representationer:

\begin{frame}[fragile]
  \inputminted[linenos,firstline=3,lastline=3]{python}{examples/myfrac_base.py}
  \dots
  \inputminted[linenos,autogobble=false,firstline=25,lastline=30]{python}{examples/myfrac_base.py}

  \begin{example}
    Vi provar att typkonvertera:
    \begin{minted}{python}
      frac = Fraction(1, 2)
      print(frac)
      print(float(frac))
    \end{minted}
  \end{example}
\end{frame}

Här ser vi två viktiga dundermetoder:
\begin{itemize}
  \item \mintinline{python}{__str__()}: returnerar en strängrepresentation för 
  utskrift (används av \mintinline{python}{print()} och 
  \mintinline{python}{str()}).
  \item \mintinline{python}{__float__()}: konverterar bråket till ett flyttal 
    (används av \mintinline{python}{float()}).
\end{itemize}

\begin{exercise}
  Hur skulle du hantera att konvertera bråket till ett heltal, exempelvis
  \mintinline{python}{int(Fraction(7, 2))}?
\end{exercise}

Ett rimligt resultat skulle vara att returnera täljaren dividerat med nämnaren, 
avrundat nedåt (heltals-division). Alltså skulle 
\mintinline{python}{int(Fraction(7, 2))} ge \mintinline{python}{3}.

För att implementera det skulle vi behöva 
\mintinline{python}{__int__()}-dundermetoden.


\section{Addition och subtraktion}

Nu kommer vi till det riktigt intressanta: att implementera aritmetiska 
operationer! Vi vill kunna skriva \mintinline{python}{frac_a + frac_b} precis 
som vi gör med heltal och flyttal.

\begin{exercise}
  Hur kan vi implementera addition av två bråk?
\end{exercise}

För att addera två bråk \(\frac{a}{b} + \frac{c}{d}\) behöver vi:
\begin{enumerate}
  \item Hitta en gemensam nämnare: \(b \cdot d\)
  \item Anpassa täljarna: \(\frac{a \cdot d}{b \cdot d} + \frac{c \cdot b}{d \cdot b}\)
  \item Addera täljarna: \(\frac{a \cdot d + c \cdot b}{b \cdot d}\)
\end{enumerate}

\subsection{Addition}

För att göra addition möjlig med \mintinline{python}{+}-operatorn behöver vi 
implementera dundermetoden \mintinline{python}{__add__()}.
Python anropar \mintinline{python}{frac_a.__add__(frac_b)} när vi kör
\mintinline{python}{frac_a + frac_b}.

\begin{frame}[fragile]
  \inputminted[linenos,firstline=3,lastline=3]{python}{examples/myfrac_add.py}
  \dots
  \inputminted[autogobble=false,linenos,firstline=31,lastline=40]{python}{examples/myfrac_add.py}

  \begin{example}
    Addera \(\frac{1}{2}\) och \(\frac{1}{3}\).
    \begin{minted}{python}
      frac_a = Fraction(1, 2)
      frac_b = Fraction(1, 3)
      print(frac_a + frac_b)
    \end{minted}
  \end{example}
\end{frame}

Nu fungerar \mintinline{python}{frac_a + frac_b}.

\begin{frame}[fragile]
  \begin{example}
    Addera \(\frac{1}{2}\) och \(1\).
    \begin{minted}{python}
      frac_a = Fraction(1, 2)
      print(frac_a + 1)
      print(1 + frac_a) # funkar ej
    \end{minted}
  \end{example}
  \begin{exercise}
    Varför fungerar inte \mintinline{python}{1 + frac_a}?
  \end{exercise}
\end{frame}

Tänk på vilken ordning Python utvärderar \mintinline{python}{1 + frac_a}. 
Vilket objekt är till vänster om \mintinline{python}{+}? Vilken metod anropas?
Python försöker först anropa \mintinline{python}{(1).__add__(frac_a)}, men 
heltalets \mintinline{python}{__add__} vet inte hur man adderar med vårt 
\mintinline{python}{Fraction}-objekt.
Vi behöver en \enquote{omvänd} additionsmetod.

\begin{frame}[fragile]
  \inputminted[linenos,firstline=3,lastline=3]{python}{examples/myfrac_add.py}
  \dots
  \inputminted[autogobble=false,linenos,firstline=41,lastline=43]{python}{examples/myfrac_radd.py}

  \begin{example}
    Prova \(1 + \frac{1}{2}\) igen.
    \begin{minted}{python}
      frac_a = Fraction(1, 2)
      print(1 + frac_a)
    \end{minted}
  \end{example}
\end{frame}

Här implementerar vi \mintinline{python}{__radd__()} (\foreignlanguage{english}{right add} 
eller \foreignlanguage{english}{reverse add}). Den anropas när det vänstra objektet 
inte kan hantera additionen.

\begin{exercise}
  Om vi redan har implementerat \mintinline{python}{__add__()}, kan vi återanvända 
  den koden i \mintinline{python}{__radd__()}? Hur?
\end{exercise}

Ja! Vi kan helt enkelt anropa \mintinline{python}{self.__add__(other)} från 
\mintinline{python}{__radd__()}, eftersom addition är kommutativ (ordningen 
spelar ingen roll).


\subsection{Negation}

Innan vi implementerar subtraktion är det användbart att kunna negera ett 
bråk, alltså göra \(\frac{1}{2}\) till \(-\frac{1}{2}\).

\begin{exercise}
  Vilken dundermetod tror du används för att implementera den unära 
  minus-operatorn (\mintinline{python}{-frac})?
\end{exercise}

Vi använder \mintinline{python}{__neg__()} för att implementera negation.

\begin{frame}[fragile]
  \inputminted[linenos,firstline=3,lastline=3]{python}{examples/myfrac_sub.py}
  \dots
  \inputminted[autogobble=false,linenos,firstline=44,lastline=45]{python}{examples/myfrac_sub.py}

  \begin{example}
    Negera \(\frac{1}{2}\).
    \begin{minted}{python}
      frac_a = Fraction(1, 2)
      print(-frac_a)
    \end{minted}
  \end{example}
\end{frame}

Negation är enkelt: vi negerar bara täljaren. \(-\frac{1}{2} = \frac{-1}{2}\).

\begin{exercise}
  Kan vi negera nämnaren istället? Vad blir skillnaden mellan \(\frac{-1}{2}\) 
  och \(\frac{1}{-2}\)?
\end{exercise}

Matematiskt är de ekvivalenta, men konventionen är att ha minustecknet i 
täljaren. Det gör representationen mer konsekvent.
Dessutom kan det uppstå problem med en negativ nämnare när vi behöver hitta 
samma bas vid aritmetiska operationer, såsom addition.


\subsection{Subtraktion}

\begin{frame}
  \begin{exercise}
    Hur kan vi implementera subtraktion av två bråk?
  \end{exercise}
\end{frame}

Nu när vi har addition och negation kan vi enkelt implementera subtraktion!
Tänk på att \(a - b = a + (-b)\).

\begin{frame}[fragile]
  \inputminted[linenos,firstline=3,lastline=3]{python}{examples/myfrac_sub.py}
  \dots
  \inputminted[autogobble=false,linenos,firstline=47,lastline=48]{python}{examples/myfrac_sub.py}

  \begin{example}
    Subtrahera \(\frac{1}{4}\) från \(\frac{1}{2}\).
    \begin{minted}{python}
      frac_a = Fraction{1, 2}
      frac_b = Fraction{1, 4}
      print(frac_a - frac_b)
    \end{minted}
  \end{example}
\end{frame}

\begin{frame}
  \begin{exercise}
    Behöver vi göra något för att subtrahera ett bråk från ett heltal, 
    exempelvis \mintinline{python}{1 - frac_a}?
    Vad?
  \end{exercise}
\end{frame}

Igen har vi samma problem som med addition: \mintinline{python}{1 - frac_a} 
fungerar inte.
Vi behöver \mintinline{python}{__rsub__()} (\foreignlanguage{english}{reverse 
subtract}), men här måste vi vara försiktiga!

\begin{frame}[fragile]
  \inputminted[linenos,firstline=3,lastline=3]{python}{examples/myfrac_sub.py}
  \dots
  \inputminted[autogobble=false,linenos,firstline=50,lastline=51]{python}{examples/myfrac_sub.py}

  \begin{example}
    Prova \(1 - \frac{1}{2}\) igen.
    \begin{minted}{python}
      frac_a = Fraction{1, 2}
      print(1 - frac_a)
    \end{minted}
  \end{example}
\end{frame}

\begin{exercise}
  Kan vi implementera \mintinline{python}{__rsub__()} på samma sätt som 
  \mintinline{python}{__radd__()}? Varför eller varför inte?
\end{exercise}

Nej! Subtraktion är \emph{inte} kommutativ:
\(1 - \frac{1}{2} \neq \frac{1}{2} - 1\).
Vi måste negera \mintinline{python}{self} och sedan addera, eller skapa ett 
nytt bråk av \mintinline{python}{other} och subtrahera 
\mintinline{python}{self}.


\section{Multiplikation}

Multiplikation av bråk är faktiskt enklare än addition!

\begin{exercise}
  Hur kan vi implementera multiplikation av två bråk, \(\frac{a}{b} \cdot 
  \frac{c}{d}\)?
\end{exercise}

Formeln är enkel: \(\frac{a}{b} \cdot \frac{c}{d} = \frac{a \cdot c}{b \cdot d}\). 
Vi multiplicerar täljare med täljare och nämnare med nämnare.

\begin{frame}[fragile]
  \inputminted[linenos,firstline=3,lastline=3]{python}{examples/myfrac_mul.py}
  \dots
  \inputminted[autogobble=false,linenos,firstline=53,lastline=59]{python}{examples/myfrac_mul.py}

  \begin{example}
    Multiplicera \(\frac{1}{2}\) och \(\frac{1}{3}\).
    \begin{minted}{python}
      frac_a = Fraction{1, 2}
      frac_b = Fraction{1, 3}
      print(frac_a * frac_b)
    \end{minted}
  \end{example}
\end{frame}

Implementationen använder \mintinline{python}{__mul__()}-dundermetoden. Notera 
att den även kan hantera multiplikation med heltal genom att skapa ett 
\mintinline{python}{Fraction}-objekt av heltalet först.

\begin{exercise}
  Vad händer om vi multiplicerar \(\frac{1}{2} \cdot 2\)? Vilket resultat 
  får vi och i vilken form?
\end{exercise}

\begin{frame}[fragile]
  \inputminted[linenos,firstline=3,lastline=3]{python}{examples/myfrac_mul.py}
  \dots
  \inputminted[autogobble=false,linenos,firstline=61,lastline=62]{python}{examples/myfrac_mul.py}

  \begin{example}
    Multiplicera \(\frac{1}{2}\) och \(2\).
    \begin{minted}{python}
      frac_a = Fraction{1, 2}
      print(frac_a * 2)
      print(2 * frac_a)
    \end{minted}
  \end{example}
\end{frame}

Vi får \(\frac{2}{2}\), vilket matematiskt är 1. Detta leder oss till nästa 
viktiga funktion: förkortning av bråk.

Precis som med addition och subtraktion behöver vi \mintinline{python}{__rmul__()} 
för att hantera när vår \mintinline{python}{Fraction} är till höger om 
multiplikationstecknet.


\section{Förkortning}

Nu har vi en fungerande bråkklass, men det finns ett problem: bråken förkortas 
inte automatiskt. Efter några operationer kan vi få \(\frac{6}{12}\) istället 
för \(\frac{1}{2}\).

\begin{frame}
  \begin{exercise}
    Hur kan vi implementera förkortning av bråk?
    Exempelvis \(\frac{2}{6}\) bör ju skrivas \(\frac{1}{3}\).
  \end{exercise}
\end{frame}

Vi kan skriva en metod som förkortar bråket.
Då kan vi anropa denna metod vid lämpliga tillfällen för att hålla bråket i 
förkortad form.

\begin{frame}
  \begin{exercise}
    Vilken matematisk operation behöver vi för att förkorta ett bråk? Vad har 
    \(\frac{2}{6} = \frac{1}{3}\) gemensamt?
    Hur kan vi beräkna denna operation i Python?
  \end{exercise}
\end{frame}

Vi behöver hitta största gemensamma delaren (\foreignlanguage{english}{Greatest Common Divisor}, 
GCD) av täljaren och nämnaren, och sedan dividera båda med den.

Python har funktionen \mintinline{python}{math.gcd(a, b)} som beräknar största 
gemensamma delaren av två tal. Vi kan använda den så här:
\mintinline{python}{gcd = math.gcd(nominator, denominator)}.

\begin{frame}
  \begin{exercise}
    När är det lämpligt att förkorta bråket?
  \end{exercise}
\end{frame}

Det finns flera möjligheter:
\begin{itemize}
  \item I konstrueraren och efter varje operation---håller bråket alltid 
    förkortat.
  \item Bara i \mintinline{python}{__str__()}---förkortar endast för visning.
  \item Som en separat metod---låter användaren välja när.
\end{itemize}

Den bästa lösningen beror på användningsfallet. För en matematisk klass är 
det vanligt att alltid hålla bråket förkortat. Vi väljer att implementera 
förkortning i konstrueraren,
för vi vill ju även ha ha bråket förkortat när vi läser ut täljaren och 
nämnaren (\mintinline{python}{Fraction.nominator} och 
\mintinline{python}{Fraction.denominator}).

\begin{exercise}
  Varför räcker det med att förkorta i konstrueraren 
  (\mintinline{python}{Fraction.__init__}) för att hålla bråket förkortat efter 
  varje operation?
\end{exercise}

Det räcker med konstrueraren eftersom varje aritmetisk operation skapar ett 
nytt \mintinline{python}{Fraction}-objekt---det vill säga: anropar 
konstrueraren.

\subsection{Implementation}

Låt oss titta på hur vi implementerar förkortning:

\begin{frame}[fragile]
  \inputminted[linenos,firstline=3,lastline=7]{python}{examples/myfrac_simplify.py}
  \dots
  \inputminted[autogobble=false,linenos,firstline=17,lastline=18]{python}{examples/myfrac_simplify.py}
  \dots
  \inputminted[autogobble=false,linenos,firstline=20,lastline=25]{python}{examples/myfrac_simplify.py}
\end{frame}

Vi behöver importera \mintinline{python}{math}-modulen för att få tillgång till 
\mintinline{python}{gcd()}-funktionen. Sedan skapar vi en (privat) metod 
\mintinline{python}{__simplify()} som gör själva förkortningen.

\begin{exercise}
  Varför är \mintinline{python}{__simplify()} en privat metod (med dubbla 
  understreck)? Behöver den inte vara publik?
\end{exercise}

Nej, det är en intern hjälpmetod som bara behövs inom klassen.
Genom att göra den privat skyddar vi klassens interna implementation.
En användare (annan programmerare) kommer aldrig att behöva den då vi alltid 
håller bråken förkortade.

Låt oss titta närmare på implementationen av \mintinline{python}{__simplify()}:

\begin{frame}[fragile]
  \inputminted[linenos,firstline=20,lastline=24]{python}{examples/myfrac_simplify.py}
\end{frame}

Metoden är enkel:
\begin{enumerate}
  \item Beräkna GCD av täljare och nämnare
  \item Dividera både täljare och nämnare med GCD (heltals-division)
\end{enumerate}

\begin{exercise}
  Varför använder vi \mintinline{python}{//=} (heltalsdivision och tilldelning) 
  istället för \mintinline{python}{/=} (vanlig division)?
\end{exercise}

Täljare och nämnare ska alltid vara heltal.
Om vi använder vanlig division får vi flyttal, vilket inte är lämpligt för en 
bråkrepresentation (rationella tal).
Eftersom att vi alltid delar med största gemensamma delaren vet vi att 
divisionen alltid går jämnt upp---det blir aldrig några decimaler.

\subsection{Testning}

\begin{frame}[fragile]
  \begin{example}
    Låt oss testa vår förkortningsfunktion:
    \only<1>{%
      \inputminted[linenos,firstline=75,lastline=92]{python}{examples/myfrac_simplify.py}
    }
    \only<2>{%
      \inputminted[linenos,firstline=75,lastline=79]{python}{examples/myfrac_simplify.py}
    }
    \dots

    \begin{onlyenv}<2>
      När vi kör detta program får vi:
      \begin{minted}{text}
        Fraction(6, 12) = 1/2
        Fraction(2, 6) = 1/3
        1/2 + 1/4 = 3/4
        1/2 * 2/3 = 1/3
        2/3 - 1/6 = 1/2
      \end{minted}
    \end{onlyenv}
  \end{example}
\end{frame}

Perfekt! Alla bråk förkortas automatiskt till sin enklaste form. 
\mintinline{python}{Fraction(6, 12)} blir \mintinline{python}{1/2}, och 
resultatet av \mintinline{python}{Fraction(1, 2) + Fraction(1, 4)} blir 
\mintinline{python}{3/4} istället för \mintinline{python}{6/8}.

\section{Sammanfattning}

Vi har nu en komplett bråkklass med:
\begin{itemize}
  \item Flexibel konstruktion av bråk
  \item Egenskaper för att läsa täljare och nämnare
  \item Aritmetiska operationer (+, -, *, med både normala och omvända varianter)
  \item Typkonvertering till sträng och flyttal
  \item Automatisk förkortning
\end{itemize}

Hela koden för klassen ser ut så här:
\begin{frame}[fragile]
  \inputminted[linenos]{python}{examples/myfrac_simplify.py}
\end{frame}

\begin{exercise}
  Vilka ytterligare funktioner skulle kunna vara användbara för en bråkklass?
\end{exercise}

Vi har en del saker som skulle kunna implementeras med fler dundermetoder:
\begin{itemize}
  \item Division av bråk, använder \mintinline{python}{__truediv__()} och 
  \mintinline{python}{__rtruediv__()}.
  \item Att skapa bråk med ett bråk i nämnaren.
  \item Hantera noll i nämnaren med undantag.
  \item Typkonvertera sträng till bråk, använder konstrueraren.
    (Då kan vi skapa bråk med \mintinline{python}{Fraction("3/4")} och således 
    använda \mintinline{python}{input_type(Fraction, "Ange ett bråktal: ")}.)
  \item Jämförelseoperatorer (\mintinline{python}{<}, \mintinline{python}{>}, 
    \mintinline{python}{==}), använder \mintinline{python}{__lt__()}, 
    \mintinline{python}{__eq__()}, etc.
  \item Upphöjning till potens, använder \mintinline{python}{__pow__()}.
  \item Absolutbelopp, använder \mintinline{python}{__abs__()}.
\end{itemize}

\begin{exercise}
  Vad händer om vi skapar ett bråk med negativ nämnare, som 
  \mintinline{python}{Fraction(1, -2)}?
  %Hur kan vi förbättra vår \mintinline{python}{__simplify()}-metod för att 
  %hantera detta?
\end{exercise}

Det är problematiskt att tillåta negativ nämnare eftersom klassen då saknar en 
entydig, kanonisk representation av ett rationellt tal.
Konsekvenserna blir:

\begin{itemize}
\item Samma tal får flera olika interna former. Exempelvis 
  \(\texttt{Fraction(1, -2)}\) och \(\texttt{Fraction(-1, 2)}\) representerar båda 
  \(-\frac{1}{2}\), medan \(\texttt{Fraction(-1, -2)}\) egentligen är \(\frac{1}{2}\).
  Om vi skriver ut dem får vi
  \mintinline{python}{print(Fraction(1, -2))} som ger \(1/-2\) och 
  \mintinline{python}{print(Fraction(-1, 2))} ger \(-1/2\), vilket ger olika 
  utskrift för samma värde.
  Detta kan leda till förvirring vid felsökning och jämförelser.
\item Förenklingen via \(\gcd\) flyttar inte tecknet till täljaren. 
  \mintinline{python}{Fraction(6, -12)} förkortas exempelvis till 
  \(\frac{1}{-2}\) i stället för \(-\frac{1}{2}\), trots att värdet är 
  detsamma.
\item Negativa nämnare \enquote{sprids} vidare i sammansatta konstruktioner:
  \mintinline{python}{Fraction(1, -2) + Fraction(1, 3)} ger \(1/-6\), medan 
  \mintinline{python}{Fraction(-1, 2) + Fraction(1, 3)} ger \(-1/6\).
  Resultatet är korrekt, men representationerna gör felsökning och jämförelser 
  svårare.
\end{itemize}

För robusthet och determinism bör man därför normalisera tecknet efter 
förenkling: se alltid till att nämnaren är strikt positiv genom att 
multiplicera både täljare och nämnare med \(-1\) när nämnaren är negativ.
Detta ger en unik, kanonisk representation för varje värde.
Vi kan lägga till kod som flyttar minustecknet till täljaren om nämnaren 
är negativ. Detta säkerställer en konsekvent representation:
\begin{minted}{python}
if self.__denominator < 0:
    self.__nominator = -self.__nominator
    self.__denominator = -self.__denominator
\end{minted}

